\documentclass[11pt]{article}
\usepackage[a4paper,margin=1in]{geometry}
\usepackage{amsmath,amssymb,mathtools,bm}
\usepackage{enumitem,booktabs,array}
\usepackage{tcolorbox}
\usepackage{hyperref}
\usepackage[protrusion=true,expansion=false]{microtype}
\hypersetup{colorlinks=true,linkcolor=blue,urlcolor=blue,citecolor=blue}
\setlist[itemize]{topsep=2pt,itemsep=2pt}
\setlist[enumerate]{topsep=2pt,itemsep=2pt}
\newtcolorbox{keybox}{colback=blue!3!white,colframe=blue!40!black,boxrule=0.4pt,arc=1mm}
\newtcolorbox{alertbox}{colback=red!3!white,colframe=red!40!black,boxrule=0.4pt,arc=1mm}
\newtcolorbox{answerbox}{colback=green!3!white,colframe=green!40!black,boxrule=0.4pt,arc=1mm}
\newtcolorbox{drillbox}{colback=yellow!5!white,colframe=orange!60!black,boxrule=0.4pt,arc=1mm}
\newcommand{\EV}{\mathbb{E}}
\newcommand{\Var}{\mathrm{Var}}
\newcommand{\Cov}{\mathrm{Cov}}
\newcommand{\blank}[1]{\underline{\hspace{#1}}}

\title{E04-04: Bernile, Bhagwat \& Rau (2017, JF) \\
\large ``What Doesn't Kill You Only Makes You More Risk-Loving: Early Life Disasters and CEO Behavior'' --- Active Recall Workbook}
\author{Empirical Corporate Finance II --- Exam Prep}
\date{\today}
\begin{document}\maketitle

\begin{keybox}
\textbf{Paper in one sentence.} BBR match CEO birth counties to county-level natural-disaster data and show a \emph{non-monotonic} relationship: CEOs who experienced \emph{moderate} natural disasters as children take more corporate risk as adults (higher leverage, volatility, risky acquisitions), while those exposed to \emph{extreme} disasters take less --- evidence that formative-experience effects on managerial behavior are U-shaped.

\textbf{Placement.} Canonical ``CEO formative-experience'' paper in behavioral corporate finance (complement to Malmendier-Nagel 2011 on depression babies and Schoar-Zuo 2017 on recession CEOs). Demonstrates non-monotonicity that simple exposure stories miss.
\end{keybox}

\section*{Round 1: Complete Walk-Through}

\subsection*{1.1 Hypothesis}
Early-life adversity may \emph{desensitize} (``what doesn't kill you makes you stronger,'' increases risk tolerance) or \emph{sensitize} (``scarring,'' decreases risk tolerance). Which dominates may depend on severity: moderate experiences desensitize, extreme ones traumatize.

\subsection*{1.2 Data and matching}
\begin{itemize}
\item ExecuComp CEOs 1992--2012; hand-collected birth counties.
\item NOAA/FEMA county-level natural disasters (hurricanes, floods, tornadoes, earthquakes, wildfires) during CEO's childhood years (age 5--15).
\item Classify exposure: none / moderate / extreme (based on fatalities).
\item Data-quality controls: historical county boundaries, migration adjustments.
\end{itemize}

\subsection*{1.3 Specification}
\[
Y_{it}=\alpha_{\text{industry}}+\delta_t+\beta_1\cdot\text{Moderate}_i+\beta_2\cdot\text{Extreme}_i+X_{it}'\gamma+\varepsilon_{it},
\]
where $Y$ includes leverage, return vol, acquisition activity, and cash holdings.

\subsection*{1.4 Headline results}
\begin{itemize}
\item Moderate-disaster CEOs: higher leverage ($+$several percentage points), higher return vol, more (and larger) acquisitions, lower cash.
\item Extreme-disaster CEOs: lower leverage, lower vol, fewer acquisitions, more cash.
\item U-shaped pattern robust to sample splits, firm controls, and matched-sample analysis.
\end{itemize}

\subsection*{1.5 Mechanism}
\begin{itemize}
\item Moderate exposure builds resilience and risk tolerance (consistent with developmental psychology on mastery).
\item Extreme exposure produces trauma; risk-avoidance dominates.
\item Effect persists through decades; consistent with formative-experience imprinting (Malmendier-Nagel).
\end{itemize}

\subsection*{1.6 Implications}
\begin{itemize}
\item CEO biographies matter for corporate policy; governance should consider this heterogeneity.
\item Provides a sharp non-monotonic test of behavioral CEO-trait channels.
\item Complements CEO-fixed-effect studies (Bertrand-Schoar 2003).
\end{itemize}

\subsection*{1.7 Caveats}
\begin{alertbox}
\begin{itemize}
\item Birth-county exposure is an imperfect proxy for personal experience.
\item Migration within childhood introduces measurement error.
\item Cannot distinguish individual-level biological/psychological channels from local-culture effects.
\end{itemize}
\end{alertbox}

\section*{Round 2: Level 1 Fill-in}
\begin{enumerate}
\item Exposure classification: none / \blank{2cm} / \blank{1.5cm}.
\item Age range used: \blank{1cm}--\blank{1cm}.
\item Moderate-disaster CEOs: \blank{1cm} risk-taking.
\item Extreme-disaster CEOs: \blank{1cm} risk-taking.
\item Pattern: \blank{1cm}-shaped exposure-risk relationship.
\end{enumerate}

\section*{Round 3: Level 2 Prompts}
\begin{enumerate}
\item State the competing hypotheses for adversity's effect on risk.
\item Describe the data construction and identification.
\item Report the U-shaped finding and its statistical robustness.
\item Relate BBR to Malmendier-Nagel (2011) and Bertrand-Schoar (2003).
\item What are the key caveats, and how might they be addressed?
\end{enumerate}

\section*{Round 4: Minimal Scaffolding}
From a blank page: (i) hypothesis, (ii) data, (iii) U-shape, (iv) mechanism, (v) caveats.

\section*{Round 5: Blank Page}
Essay: early-life disasters, CEO risk preferences, and corporate policy.

\vspace{6pt}
\begin{drillbox}
\textbf{Speed Drill.} (1) Exposure categories. (2) Childhood age range. (3) Moderate-exposure effect. (4) Extreme-exposure effect. (5) Pattern shape.
\end{drillbox}

\section*{Quick Reference \& Answer Key}
\begin{answerbox}
\begin{itemize}
\item \textbf{Design.} County-level childhood disaster exposure.
\item \textbf{Moderate.} $\uparrow$ risk-taking.
\item \textbf{Extreme.} $\downarrow$ risk-taking.
\item \textbf{Pattern.} U-shaped in severity.
\end{itemize}
\end{answerbox}

\begin{keybox}
\textbf{30-second exam pitch.} Bernile, Bhagwat \& Rau (2017) test whether early-life adversity affects adult CEOs' corporate risk-taking. They hand-match ExecuComp CEOs (1992--2012) to their birth counties and merge with NOAA/FEMA county-level natural-disaster data for the CEOs' childhood years. Classifying exposure as none, moderate (disaster present), or extreme (disasters with significant fatalities), they find a U-shaped pattern: moderate-disaster CEOs take more corporate risk as adults (higher leverage, return vol, more risky M\&A, lower cash) while extreme-disaster CEOs take less. The effect is robust across specifications and sample splits. Interpretation: moderate adversity builds resilience and risk tolerance; extreme adversity traumatizes. The paper is a canonical CEO-formative-experience study in behavioral corporate finance, complementing Malmendier-Nagel's depression-babies and Bertrand-Schoar's manager-fixed-effects work, and cautioning that simple ``adversity increases risk tolerance'' stories miss important non-monotonicity.
\end{keybox}

\end{document}
